
% --- SLIDE 3: SEEDING ---
\begin{frame}[fragile]{Seeding: Every Library, Every Time}

    \begin{lstlisting}[language=Python, basicstyle=\ttfamily\scriptsize]
# PYTHONHASHSEED must be set BEFORE the interpreter starts, so it goes in the
# launcher, not in the script:   export PYTHONHASHSEED=42
import random, numpy as np, torch

def set_seed(seed: int = 42):
    random.seed(seed)          # Python stdlib
    np.random.seed(seed)       # NumPy legacy global RNG
    torch.manual_seed(seed)    # CPU *and* all CUDA devices

set_seed(42)

# scikit-learn: pass random_state explicitly -- never rely on the global seed
train_test_split(X, y, test_size=0.2, random_state=42, stratify=y)
RandomForestClassifier(n_estimators=500, random_state=42)
    \end{lstlisting}

    \vspace{0.3em}
    \begin{itemize}
        \small
        \item \textbf{Never rely on the global NumPy seed for scikit-learn.} Anything left at
              \texttt{random\_state=None} falls back to NumPy's \emph{global} RNG, so its result depends on
              whatever consumed random numbers first. Pass \texttt{random\_state} explicitly, everywhere.
        \item Any subprocess gets a fresh interpreter and therefore a fresh RNG state --- PyTorch
              \texttt{DataLoader} workers, \texttt{joblib} parallel jobs, Spark executors. Each needs its
              own seeding hook.
        \item \textbf{Log the seed as a tracked parameter.} A seed you cannot look up later is not a seed.
    \end{itemize}

\end{frame}
